%قبل از شروع مسئله ترجیحاً چیزی ننویسید. اما اگر خواستید بنویسید، در قالب لیستی (فقرات و شمارش) بنویسید.
\شروع‌مسئله[33] \متن‌سیاه{سؤال‌های جواب کوتاه.} در سؤال‌های این قسمت معمولاً یک جواب کوتاه (۱ نمره) و یک توضیح کوتاه حدوداً بین یک تا دو خط (۲ نمره) مدنظر است؛ مگر این که صورت آن قسمت به صراحت از شما خواستۀ دیگری داشته باشد.
% صورت اصلی مسئله را در این‌جا بنویسید.
\شروع{زیرمسائل}
\زیرمسئله[3]
آیا مرتبۀ زمانی یک الگوریتم برابر با تعداد مراحل آن الگوریتم در مدل محاسباتی‌اش است؟

\پاسخ{
	خیر، زیرا ممکن است زمان اجرای یک مرحله در مدل محاسبات $O(1)$ نباشد (مثل ماشین پایتون که در کلاس گفته شد). این سؤال در کوییز ۱ نیز مطرح شده بود.
	
}
\زیرمسئله[3]
آیا الگوریتمی وجود دارد که $k$ لیست مرتب، که هر لیست به طول $n$ است، را در زمان $O(nk)$ (در یک لیست مرتب به طول $nk$) ادغام کند؟

\پاسخ{
خیر، در حالت حدی فرض کنید که لیست‌ها تک عنصره باشند ($n=1$). در این حالت ادغام $k$ لیست تک‌عنصره، همانند مرتب‌سازی $k$ عدد می‌شود که نمی‌توان (در مدل مقایسه) مرتب‌سازی را در $O(k)$ انجام داد.
	
}
\زیرمسئله[3]
آرایۀ $A$ یک هرم بیشینه را نشان می‌دهد که یک عنصر آن تغییر کرده است و رأس مربوط به آن دیگر خاصیت هرم بیشینۀ رأسی را ندارد. ساختار درختی این هرم را ترسیم کنید و رأس گفته شده را در آن مشخص کنید.
\begin{flushleft}$
	A=[18, 15, 12, 9, 14, 13, 8, 7, 6]$
\end{flushleft}
\پاسخ{
	عنصر ۱۲ خاصیت هرم بیشینه را ندارد (زیرا ۱۳ فرزند آن است). توجه کنید ۱۳ خاصیت هرم بیشینۀ رأسی را دارد زیرا از فرزندان خود بزرگتر است (زیرا فرزندی ندارد).

}

\نمره‌دهی{
	اگر رأس ۱۳ را به جای ۱۲ نوشته باشند، ۱.۵ نمره تعلق می‌گیرد.
	
}

\زیرمسئله[3]
آیا می‌توان در زمان خطی یک هرم بیشینه را به یک هرم کمینه تبدیل کرد؟

\پاسخ{
	بله، هر آرایۀ  دلخواه از اعداد را می‌توان در زمان خطی به هرم کمینه تبدیل کرد.

}

\زیرمسئله[3]
آیا مرتب‌سازی هرمی روی یک آرایۀ مرتب در زمان $O(n)$ انجام می‌شود؟

\پاسخ{
	خیر، زیرا در مرتب‌سازی هرمی در $n$ مرحله عنصر بیشینه استخراج می‌شود که هر کدام از این استخراج‌ها $\Theta(\log n)$ زمان می‌گیرد.
	
}

\زیرمسئله[3]
در یک درخت دودویی جست‌وجو، تعدادی عنصر با کلید یکسان $a$ درج شده‌اند. آیا ممکن است یک رأس با کلید $a$ در این درخت وجود داشته باشد که کلید‌های والد و فرزندان آن $a$ نباشند؟ (اگر جواب خیر است دلیل آن را توضیح دهید و اگر جواب بله است یک مثال بزنید)

\پاسخ{
بله. یک مثال سه رأسی برای آن وجود دارد که کلید ریشه ۲، فرزند چپ ریشه دارای کلید ۱ باشد و خود این رأس فرزند راستی داشته باشد که دارای کلید ۲ باشد.
	
}

%\زیرمسئله[4]
%این جمله درست است یا غلط؟ اگر عنصر با کلید $a$ در یک درخت AVL برگ نباشد، عنصر با کوچکترین  کلید بزرگ‌تر از $a$ یک برگ است.

%\زیرمسئله[4]
%ترتیبی از اعداد ۱ تا ۷ را بیان کنید که با درج آن‌ها در یک درخت AVL کمترین تعداد دوران‌ها را داشته باشد؟ در ابتدا درخت AVL تهی است.

\زیرمسئله[3]
آیا یک درخت دودویی جست‌وجوی غیرمتوازن را می‌توان به یک درخت AVL با حداکثر $O(\log n)$ دوران تبدیل کرد؟

\پاسخ{
	خیر، زیرا هر دوران ارتفاع را حداکثر یک واحد کم می‌کند و مثال‌هایی از درخت‌های دودویی جست‌وجوی غیرمتوازن وجود دارد که ارتفاع آن‌ها 
	$\Theta(n)$
	 است.
	
}

\نمره‌دهی{
\شروع{فقرات}
\فقره کسانی که به $\mathcal{O}(n)$ بودن اشاره کرده‌اند یک نمره گرفته‌اند.
\فقره اشاره به اینکه با هر دوران می‌توان ارتفاع را حداکثر یک واحد کاهش داد یک نمره دارد و کسانی که به این موضوع اشاره نکرده‌اند یک نمره از دست داده‌اند.
\پایان{فقرات}

}

%\زیرمسئله[3]
%یک درخت قرمز-سیاه مثال بزنید که AVL نباشد.

\زیرمسئله[3]
آیا یک الگوریتم مونت‌کارلو همیشه جواب درست را برمی‌گرداند؟

\پاسخ{
	خیر، الگوریتم‌های مونت‌کارلو الگوریتم‌های احتمالاً درست هستند که درستی خروجی الگوریتم احتمالی است ولی همیشه سریع هستند. در کوییز سوم مشابه این سؤال وجود دارد.
	
}

\زیرمسئله[3]
یک درخت قرمز-سیاه با $n^{10}$ رأس داریم. یک جست‌وجو در این درخت (به صورت مجانبی) چه زمانی دارد؟

\پاسخ{
ارتفاع درخت (به خاطر توازن)
$O(\log (n^{10})) = O(10 \cdot\log n) = O(\log n)$
است و جست‌وجو در همین مرتبه انجام می‌شود.
	
}

\نمره‌دهی{
\شروع{فقرات}
\فقره کسانی که به $\mathcal{O}(\log n)$ بودن جست‌و‌جو در درخت قرمز-سیاه اشاره کرده‌اند یک نمره گرفته‌اند.
\فقره کسانی که $n^{10}$ را در رابطه بالا جایگذاری نکرده‌اند، یا جواب را به صورت ساده نشده $\mathcal{O}(10 \log n)$  
ارائه داده‌اند یک نمره از دست داده‌اند.

\پایان{فقرات}

}

\زیرمسئله[3]
آیا ۷ عنصر را می‌توان با ۱۲ مقایسه مرتب کرد؟

\پاسخ{
	خیر، زیرا بنابر استدلال درخت تصمیم به
	$\lceil\log_2(7!)\rceil = \lceil\log_2(5040)\rceil = 13$
	مقایسه نیاز داریم.
	
}

\نمره‌دهی{
\شروع{فقرات}
\فقره کسانی که به درخت تصمیم اشاره کرده‌اند یک نمره دریافت کرده‌اند.
\فقره کسانی که در محاسبه اشتباه کرده‌اند یک نمره از دست داده‌اند.
\فقره کسانی که گفته‌اند مرتب‌سازی در بهترین حالت در $\mathcal{O}(n\log n)$
است و محاسبه را نیز با 
$n \log n$
انجام داده‌اند یک نمره از دست داده‌اند.

\پایان{فقرات}

}


\زیرمسئله[3]
رأس‌های درخت زیر را با رنگ‌های قرمز و سیاه طوری رنگ کنید که خاصیت‌های درخت قرمز-سیاه را داشته باشد. رنگ برگ‌ها را نیز مشخص کنید.

\پاسخ{
	رنگ رأس‌های ۱۵، ۷ و ۱۳ قرمز و رنگ بقیه سیاه است. این سؤال در کوییز دوم بوده است.
	
}

%\begin{figure}[h!]%[5]{L}{0.4\textwidth}
	%	\vspace{-2cm}
	\begin{center}
		\begin{tikzpicture}[-,>=stealth',level/.style={sibling distance = 5cm/#1,
			level distance = 1cm}] 
		\node [arn_s] {5}
		child{ node [arn_s] {1} 
			child{ node [arn_x] {}}
			child{ node [arn_x] {}}
		}
		child{ node [arn_s] {15}
			child{ node [arn_s] {10} 
				child{ node [arn_s] {7}
					child{ node [arn_x] {}}
					child{ node [arn_x] {}}
				}
				child{ node [arn_s] {13}
					child{ node [arn_x] {}}
					child{ node [arn_x] {}}
				}
			}
			child{ node [arn_s] {20}
				child{ node [arn_x] {}}
				child{ node [arn_x] {}}
			}
		}
		; 
		\end{tikzpicture}
	\end{center}
%	\caption{درخت قرمز-سیاه}
%	\label{شکل:۷و۸}
%\end{figure}

\نمره‌دهی{
\شروع{فقرات}
\فقره در این قسمت برای رنگ برگ‌ها یک نمره در نظر گرفته شده است.
\پایان{فقرات}

}


\پایان{زیرمسائل}



%اگر مسئله چند قسمت داشت به این صورت عمل کنید:
%\شروع{زیرمسائل}
%\زیرمسئله[0] % مسئلۀ ۱آ
%\زیرمسئله[0] % مسئلۀ ۱ب
%\زیرمسئله[0] % مسئلۀ ۱پ
%\زیرمسئله[0] % مسئلۀ ۱ت
%\پایان{زیرمسائل}




%%%%%%%%%%%%%%%%%%%%%%%%%%%%%%%%%%%%%%%%%%%%%%%%%%%%%
% برخی دستورات مفید برای تمرین‌ها. %
%%%%%%%%%%%%%%%%%%%%%%%%%%%%%%%%%%%%%%%%%%%%%%%%%%%%%

% فرمول‌های نمونه:
%$f(x) = \Theta(x)$
%$T(x, y) \leq \log(x) + 2^y + \binom{2n}{n}$
% {\tt code\_function}


% لیست موارد بدون شماره‌گذاری:
%\شروع{فقرات}
%    \فقره مورد اول 
%        \شروع{فقرات}
%            \فقره مورد اول داخلی 
%                \شروع{فقرات}
%                    \فقره اولین مورد خیلی داخلی 
%                    \فقره دومین مورد خیلی داخلی 
%                \پایان{فقرات}
%            \فقره مورد دوم داخلی 
%        \پایان{فقرات}
%    \فقره مورد دوم 
%\پایان{فقرات}

% لیست موارد شماره‌گذاری شده:
%\شروع{شمارش}
%    \فقره یک مورد 
%    \فقره یک مورد دیگر 
%    \فقره باز هم یک مورد دیگر
%\پایان{شمارش}

% برای نوشتن یک فرمول چند خطی که یک شماره بخوره. برای حالت‌های دیگر به سایت زیر مراجعه کنید:
% https://www.overleaf.com/learn/latex/Aligning_equations_with_amsmath
%\begin{align} 
%    \begin{split}
%        (x+y)^3 &= (x+y)^2(x+y) \\
%                &= (x^2+2xy+y^2)(x+y) \\
%                &= (x^3+2x^2y+xy^2) + (x^2y+2xy^2+y^3) \\
%                &= x^3+3x^2y+3xy^2+y^3
%    \end{split}                                 
%\end{align}

% نوشتن ماتریس و آرایه دوبعدی
%\begin{align}
%    A = 
%    \begin{bmatrix}
%        A_{11} & A_{21} \\
%        A_{21} & A_{22}
%    \end{bmatrix}
%\end{align}

% اضافه کردن تصویر با شرح. برای اضافه کردن تصویر بدون شرح، فقط دستور درج‌تصویر را استفاده کنید.
%\شروع{شکل}[h!]
%\تنظیم‌ازوسط
%  \درج‌تصویر[پهنا=0.5\پهنای‌متن]{img.jpg}
%\شرح{شرح تصویر}
%\برچسب{تصویر:۱}
%\پایان{شکل}
%\پایان{زیرمسائل}

%کد زیر را ببینید.
%%\lstinputlisting{source_filename.py}
%\begin{latin}
%\begin{lstlisting}
%# Import the modules
%import sys
%import random
%
%ans = True
%
%while ans:
%question = raw_input("Ask the magic 8 ball a question: (press enter to quit) ")
%
%answers = random.randint(1,8)
%
%if question == "":
%sys.exit()
%
%elif answers == 1:
%print "It is certain"
%
%elif answers == 2:
%print "Outlook good"
%
%elif answers == 3:
%print "You may rely on it"
%
%elif answers == 4:
%print "Ask again later"
%
%elif answers == 5:
%print "Concentrate and ask again"
%
%elif answers == 6:
%print "Reply hazy, try again"
%
%elif answers == 7:
%print "My reply is no"
%
%elif answers == 8:
%print "My sources say no"
%\end{lstlisting}
%\end{latin}
